\section[几何 Satake（终）]{几何 Satake（终）（2 月 5 日）}

\srcpage{133}

\subsection{回顾}

设 $E$ 是非阿基米德局部域，剩余域为 $\Fq\subset\ol{\Fq}$，$\breveE$ 照常。设 $G/E$ 是任意约化群。

我们定义了 Beilinson--Drinfeld Grassmann 簇 $\Gr^I_G\to(\Div^1_X)^I$，其中
\[
\Div^1_X=(\Spa\breveE)^\diam/\varphi^{\ZZ}
\]
是 $\Perf_{\Fq}$ 上的 $v$-层。群 $(L^+G)^I$ 作用在 $\Gr^I_G$ 上；我们还定义了 $(\Div^1_X)^I$ 上的
\[
\mathrm{Hk}^I_G=(L^+G)^I\backslash\Gr^I_G .
\]
此外，我们定义了反常层与 ULA 层的概念。

\begin{definition}[Satake 范畴]
设 $\Lambda$ 是环，使对某个与 $p$ 互素的 $n$ 有 $n\Lambda=0$。则
\[
\mathrm{Sat}^I_G(\Lambda):=\Perv^{ULA}_{flat}\bigl(\mathrm{Hk}^I_G,\Lambda\bigr),
\]
带有纤维函子\footnote{这里 $\LocSys$ 与 $\Rep$ 都指有限投射 $\Lambda$-模上的相应范畴。}
\[
F^I=\bigoplus_i R^i\pi_{G*}\colon\mathrm{Sat}^I_G(\Lambda)\longrightarrow
\LocSys\bigl((\Div^1_X)^I,\Lambda\bigr)\cong\Rep_{W^I_E}(\Lambda).
\]
\end{definition}

\begin{remark}
对 $F^I$ 有一个「更好」的描述，至少当 $G$ 有 Borel 子群 $B$ 时如此。（一般情形用下降，从这一情形逐步搭起来。）我们有常数项函子
\[
\mathrm{CT}_B=Rp_!q^*[\deg]\colon\mathrm{Sat}^I_G(\Lambda)\to\mathrm{Sat}^I_T(\Lambda),
\]
其中 $p,q$ 如下面的交换图所示：
\[
\begin{tikzcd}
& \Gr^I_B \arrow[dl, "q"'] \arrow[dr, "p"] & \\
\Gr^I_G & & \Gr^I_T
\end{tikzcd}
\]
这个函子其实与 $B$ 无关。事实上，让 $B$ 变动，我们得到一个到
$\LocSys\bigl((\Div^1_X)^I\times\mathrm{Fl}^\diam_G,\Lambda\bigr)$ 的函子；但旗簇
$\mathrm{Fl}^\diam_G$ 单连通，所以这样的局部系统都从底空间
$\LocSys\bigl((\Div^1_X)^I,\Lambda\bigr)$ 拉回。特别地，$\mathrm{CT}_B$ 下降为所有 $G$ 上的函子。由双曲局部化，有
\[
\pi_{T*}\mathrm{CT}_B\cong\bigoplus_i R^i\pi_{G*}=F^I .
\]
\end{remark}

融合积使 $\mathrm{Sat}^I_G$ 成为对称幺半范畴，而
$F^I\colon\mathrm{Sat}^I_G\to\Rep_{W^I_E}(\Lambda)$ 是对称幺半纤维函子。再多做一点工作可见，存在 Hopf 代数 $\mathcal{H}^I\in\Ind\Rep_{W^I_E}(\Lambda)$，使得
\[
\mathrm{Sat}^I_G\cong\mathrm{CoMod}_{\mathcal{H}^I}\bigl(\Rep_{W^I_E}(\Lambda)\bigr).
\]
「Künneth 公式」蕴含 $\mathcal{H}^I\cong\bigotimes_{i\in I}\mathcal{H}^{\{i\}}$，因此只需确定 $\mathcal{H}^{\{1\}}$。它等价于一个仿射平坦群概形 $\check G_\Lambda/\Lambda$ 连同连续的 $W_E$-作用。$\mathcal{H}^{\{1\}}$ 的构造与 $\Lambda$ 的基变换相容，所以可以假设 $\Lambda=\ZZ/\ell^n\ZZ$。

\begin{theorem}
设 $\hat G$ 是 Langlands 对偶群（下文解释）。存在典范同构
\[
\check G_{\ZZ/\ell^n\ZZ}\cong\hat G_{\ZZ/\ell^n\ZZ},
\]
若 $\hat G$ 上的 $W_E$-作用带有分圆扭转（下文解释），则它是 $W_E$-等变的。
\end{theorem}

\begin{remark}
记号上，我们把由几何 Satake 构造出的群 $\check G$ 与抽象构造的对偶群 $\hat G$ 区分开。
\end{remark}

\srcpage{134}

\begin{corollary}
存在典范等价
\[
\mathrm{Sat}^I_G(\Lambda)\cong\Rep_{\Rep_{W^I_E}(\Lambda)}\bigl(\hat G^I_\Lambda\bigr),
\]
它与到 $\Rep_{W^I_E}(\Lambda)$ 的各自遗忘函子交换。换言之，只要后者被适当地定义，就有
\[
\mathrm{Sat}^I_G(\Lambda)\cong\Rep\bigl((\hat G_\Lambda\rtimes W_E)^I\bigr).
\]
\end{corollary}

\begin{remark}
在非代数闭域上，完整的 $L$-群也能从几何 Satake 中产生；这一点最早由 Timo Richarz 与 Xinwen Zhu 观察到。
\end{remark}

\subsection{对偶群}

设 $G/E$（或任意域上的 $G$）。

\subsubsection{泛 Cartan}

$G$ 有一个「泛 Cartan」$T$。可以这样设想：存在旗簇 $\mathrm{Fl}_G/E$，它参数化 Borel 子群 $B\subset G$，而每个 $B$ 都有环面商 $T$。这个族对应于 $\mathrm{Fl}$ 上的一个 $\ZZ$-局部系统 $X^*(T)$。（概形上的环面与 $\ZZ$-局部系统相对应，把环面送到其特征群；这是一般事实。）

现在，几何旗簇单连通，因此 $\mathrm{Fl}_G$ 上任何 $\ZZ$-局部系统都从底空间拉回，此处底空间是 $(\Spec E)_{\et}$。于是 $T$ 的族唯一地由某个环面 $T/E$ 拉回得到。

特别地，$\Gal(\ol E/E)$ 作用在 $X^*(T_{\ol E})$ 上。于是从 $G/E$ 得到一个有限自由 $\ZZ$-模 $X^*$，连同 $\Gal(\ol E/E)$-作用。

\subsubsection{根数据}

但事实上我们得到的东西略多一点。事实上，泛 Cartan 与任一 Borel 的商等同，因此它带有典范的支配余特征 $X^*_+$ 的概念。这完全是典范的，因而在 $\Gal(\ol E/E)$ 作用下稳定。此外还有 Weyl 群 $W$ 在 $X^*$ 上的作用。

既然有支配余特征，也就有支配 Weyl 房，于是我们还得到一组单反射 $S\subset W$，它在 $\Gal(\ol E/E)$-作用下保持不变。对任意单反射 $s\in S$，有：
\begin{itemize}
\item 单根 $\alpha_s\in X^*$。给定一个 Borel $B$，得到根空间 $U_{\alpha_s}\subseteq G_{\ol E}$。
\item 单余根 $\alpha_s^\vee\in X_*:=(X^*)^\vee$。给定一个 Borel $B$，可以把它看作来自某个对应于 $s$ 的映射 $\mathrm{SL}_2\to G_{\ol E}$。
\end{itemize}
所以，不依赖任何选择，我们有根数据 $(X^*,X_*,\Phi,\Phi^\vee,X^*_+,X^+_*)$，连同 $\Gal(\ol E/E)$ 的作用。

\subsubsection{Langlands 对偶群}

现在，Langlands 观察到：把 $X^*$ 与 $X_*$ 交换，同样给出根数据。从（约化群）到（根数据）的函子有典范分裂，由 Chevalley 群概形给出。我们对 $(X_*,X^*,\Phi^\vee,\Phi,X^+_*,X^*_+)$ 取对应的 Chevalley 群概形，记作 $\hat G/\ZZ$。由于 Chevalley 群概形是函子式构造，它继承 $\Gal(\ol E/E)$ 的作用。这就是 $G$ 的「对偶群」。

注意 $\hat G$ 带有典范的 $\hat T\subset\hat B\subset\hat G$，它们特别地在 $\Gal(\ol E/E)$-作用下稳定。我们还有典范等同 $X^*(\hat T)=X_*$；并且作为 Chevalley 构造的一部分，对任意单反射 $s\in S$ 有同构
\[
\psi_s\colon\Lie\hat U_{\check\alpha_s}\cong\ZZ,
\]
它们是 $\Gal(\ol E/E)$-不变的。

\subsection{分圆扭转}

现在我们改在 $\ZZ_\ell$ 或 $\ZZ/\ell^n\ZZ$ 上工作，而不再用 $\ZZ$；同时把 $\psi_s$ 换成同构 $\psi'_s\colon\Lie\hat U_{\check\alpha_s}\cong\ZZ_\ell(1)$。这是 $\Gal(\ol E/E)$ 的 Tate 扭转表示：抽象群还是同一个，但我们改变了 $W_E$-作用。

为证明定理 25.3，需要找到 $W_E$-稳定的 $\check T\subset\check B\subset\check G$，使得：
\begin{itemize}
\item $X^*(\check T)=X_*$。
\item $\check G$ 是正确类型的约化群。
\srcpage{135}
\item 有同构 $\Lie\check U_{\check\alpha_s}\cong\ZZ/\ell^n\ZZ(1)$。
\end{itemize}

\subsection{定理 25.3 的证明}

由于我们会典范地给出全部所需结构，其实不必追踪 Galois 等变性。只要构造不依赖分裂，就可以基变换并假设 $G$ 分裂。不依赖性如上所证：用旗簇上的局部系统来陈述，再利用它单连通。这个论证已经出现过多次，今后不再重复。

\subsubsection{环面的情形}

有
\[
\mathrm{Sat}_T\cong\Perv^{ULA}_{flat}(\mathrm{Hk}_T).
\]
由于
\[
\mathrm{Hk}_T=\bigsqcup_{X_*(T)}[\Div^1_X/L^+T],
\]
直接看出
\[
\Perv^{ULA}_{flat}(\mathrm{Hk}_T)=\bigoplus_{X_*(T)}\Rep_{W_E}(\Lambda).
\]
而范畴 $\bigoplus_{X_*(T)}\Rep_{W_E}(\Lambda)$ 可以典范地等同于 $\Rep(\check T)$，其中 $\check T$ 是满足 $X^*(\check T)=X_*(T)$ 的环面。于是有典范等同
\[
\mathrm{Sat}_T\cong\Rep(\check T).
\]

\subsubsection{}

函子 $\mathrm{CT}_B\colon\mathrm{Sat}_G\to\mathrm{Sat}_T$ 是对称幺半的，并与各自的纤维函子交换。它诱导映射 $\check T\to\check G$。要检验它是单射，可以考察标准对象 $j_{\mu!}\Lambda$，计算纤维函子 $F$，并观察到最高权空间是一维的。

此外，纤维函子 $F=\bigoplus_i R^i\pi_{G*}$ 上的上同调 $\ZZ$-分次给出映射 $\GG_m\to\check G$；可以检验它典范地分解为
\[
\begin{tikzcd}
\check T \arrow[rr] & & \check G \\
& \GG_m \arrow[ul] \arrow[ur] &
\end{tikzcd}
\]
（这归结于：上同调分次与 $\mathrm{CT}$ 定义中的次数平移一致。）这个 $\GG_m$ 定义了吸引的「抛物」$\check T\subset\check B\subset\check G$（但先验地我们不知道 $\check B$ 是抛物子群）。

\subsubsection{秩 1 的群}

现在设 $G$ 的秩为 $1$，于是 $G_{\mathrm{ad}}\cong\mathrm{PGL}_2$。映射 $G\twoheadrightarrow G_{\mathrm{ad}}$ 诱导
\[
\Gr_G=\Gr_{G_{\mathrm{ad}}}\times_{\ZZ/2\ZZ}\pi_1(G).
\]
（特别地，在连通分支上 $\Gr_G$ 与 $\Gr_{G_{\mathrm{ad}}}$ 没有差别。）因此
\[
\check G\cong\check G_{\mathrm{ad}}\times_{\mu_2}\check Z,
\]
其中 $\check Z$ 是满足 $X^*(\check Z)=\pi_1(G)$ 的环面。

\begin{remark}
事实上，同样的技巧可以把一般情形化归到 $G_{\mathrm{ad}}$。
\end{remark}

总之，我们已化归到分析 $G=\mathrm{PGL}_2$。它有极小的余特征 $\mu$，且 $\Gr_{\mathrm{PGL}_2,\mu}\cong(\PP^1)^\diam$。于是 $j_{\mu*}\Lambda[1]\in\mathrm{Sat}_G(\Lambda)=\Rep(\check G)$。所以
\[
F(A)=H^0(\PP^1,\Lambda)\oplus H^2(\PP^1,\Lambda)=\Lambda\oplus\Lambda(-1).
\]
有一个重言映射 $\check G\to\Aut(F(A))=\mathrm{GL}_2/\Lambda$。

\srcpage{136}

我们知道，在 $\Ql$ 与 $\Fl$ 上，
\[
\mathrm{Irr}\Rep(\check G)\leftrightarrow X^+_*=\ZZ_{\ge0},
\]
其中 $\lambda\in X^+_*$ 对应 $\mathrm{IC}_\lambda$。这里，我们定义
\[
\mathrm{Sat}_G(\ZZ_\ell):=\varprojlim_n\mathrm{Sat}_G(\ZZ/\ell^n\ZZ)
\]
以及
\[
\mathrm{Sat}_G(\Ql):=\mathrm{Sat}_G(\ZZ_\ell)[1/\ell].
\]
我们还知道 $\mathrm{Sat}_G(\Ql)$ 是半单的（所有对象都是 $\mathrm{IC}_\lambda$ 的直和），这通过与 Witt 向量仿射 Grassmann 簇作比较得到。

这种半单性蕴含 $\check G_{\Ql}$ 是约化的（包括连通性——通过检验表示范畴没有真子张量范畴，并依据不可约对象的分类得到），且秩为 $1$（再次依据不可约对象的分类）。这意味着从 $\check G_{\Ql}$ 出发的映射在 $\Ql$ 上必须经过 $\check G\to\mathrm{SL}_2\subset\mathrm{GL}_2$ 分解，因此在 $\ZZ_\ell$ 上也如此。考察环面上发生的情况——它由 $F(A)$ 上的分次决定——
\[
\begin{tikzcd}
\check G \arrow[rr] & & \mathrm{SL}_2 \\
& \GG_m \arrow[ul] \arrow[ur] &
\end{tikzcd}
\]
可见 $\check G\to\mathrm{SL}_2$ 在 $\Ql$ 上是同构。

在 $\Fl$ 上，$\check G_{\Fl}\to\mathrm{SL}_{2,\Fl}$ 必须是满射；否则其像会是某个包含该环面的子群，而在此情形下仅有的其他可能是该环面、Borel 子群或环面的正规化子——这些可能性都会导致过多的不可约表示。

因此 $\ZZ_\ell$ 上的映射 $\cO(\mathrm{SL}_2)\to\cO(\check G)$ 在特征 $0$ 下是同构，两者都在 $\ZZ_\ell$ 上平坦，且它模 $\ell$ 是单射。于是由一般代数它是同构。（需要检验映射是满射：任何元素在乘以 $\ell^n$ 后都落在像中，但模 $\ell$ 的归约是 $0$……）

所以 $\check G\xrightarrow{\sim}\mathrm{SL}_2(F(A))=\mathrm{SL}_2(\Lambda\oplus\Lambda(-1))$。这就是在根群中看出分圆扭转的方式。

\subsection{一般 $G$}

对任意抛物 $P\twoheadrightarrow M$，$\mathrm{CT}_P$ 存在。这定义了对称幺半函子
\[
\begin{tikzcd}
\mathrm{Sat}_G \arrow[r, "\mathrm{CT}_P"] \arrow[dr] & \mathrm{Sat}_M \arrow[r, "\mathrm{CT}_{B_M}"] \arrow[d] & \mathrm{Sat}_T \arrow[dl] \\
& \Rep_{W_E} &
\end{tikzcd}
\]
对极小抛物使用它，于是 $M$ 的秩为 $1$。此时我们知道
\[
\begin{tikzcd}
\check T \arrow[r] \arrow[d, "\cong"] & \check M \arrow[r] \arrow[d, "\cong"] & \check G \\
\hat T & \hat M &
\end{tikzcd}
\]
特别地，对任意单反射 $s$ 得到
\[
\begin{tikzcd}
\GG_m \arrow[r, hook] & \mathrm{SL}_2 \arrow[r, hook, "\alpha_s\in X^*=X_*(\check T)"] & \check G
\end{tikzcd}
\]

\srcpage{137}

此时我们知道 $\check G_{\Ql}$ 连通且约化，并且至少有与 $\hat G$ 一样多的根与余根。我们需要检验它不会更多。为此使用组合论证：考察 $F(\mathrm{IC}_\lambda)$ 中出现的权，即 $\check G$ 的对应不可约表示中出现的权。这个数目已经与人们对 $\hat G$ 的最高权 $\lambda$ 表示的预期相同。这说明 $\check G_{\Ql}\cong\hat G_{\Ql}$ 是典范的，因为单根子群已被上面关于 $\alpha_s$ 的映射钉住。

我们还得到了整系数的 $\check M\to\check G$。特别地，$\check G(\breve{\ZZ}_\ell)\subset\check G(\breve{\QQ}_\ell)\cong\hat G(\breve{\QQ}_\ell)$ 含有 $\check M(\breve{\ZZ}_\ell)\subset\hat G(\breve{\QQ}_\ell)$，对任意极小 Levi 都如此。这些生成整个 $\hat G(\breve{\ZZ}_\ell)$，所以 $\hat G(\breve{\ZZ}_\ell)$ 必包含于 $\check G(\breve{\ZZ}_\ell)\subset\hat G(\breve{\QQ}_\ell)$。但 $\check G(\breve{\ZZ}_\ell)$ 有界，而 $\hat G(\breve{\ZZ}_\ell)$ 是超特殊极大紧子群，因此必须有 $\hat G(\breve{\ZZ}_\ell)=\check G(\breve{\ZZ}_\ell)$。再多做一点工作可见，作为 $\check G_{\Ql}\cong\hat G_{\Ql}$ 的整模型，有 $\check G=\hat G$。具体地，我们用：

\begin{lemma}[Prasad--Yu]
设 $H$ 是 $\ZZ_\ell$ 上的约化群，$H'$ 是 $\ZZ_\ell$ 上有限型的仿射平坦群概形，$\rho\colon H\to H'$ 是态射，且在一般纤维上是闭浸入。假设 $\ell\ne2$，或 $H_{\Ql}$ 的任何几乎单因子都不同构于 $\mathrm{SO}_{2n+1}$（例如当 $H$ 的导出群单连通时满足此条件）。则 $\rho$ 是闭嵌入。
\end{lemma}

如何应用？可以假设 $G$ 是伴随的，于是 $\hat G$ 单连通。（这意味着我们不遇到引理中禁止的情形。）取表示 $\check G\to\mathrm{GL}_N$，它在一般纤维上是闭浸入。（这里要牢记一个敌人：我们事先不知道 $\check G$ 是有限型的。嵌入到 $\mathrm{GL}_N$ 消除了这一顾虑。）于是得到 $\hat G_{\Ql}\cong\check G_{\Ql}\hookrightarrow\mathrm{GL}_{N,\Ql}$ 与 $\hat G(\breve{\ZZ}_\ell)\cong\check G(\breve{\ZZ}_\ell)\subset\mathrm{GL}_N(\breve{\ZZ}_\ell)$。由于 $\hat G$ 与 $\mathrm{GL}_N$ 是光滑群概形，可以在点上检验整结构，因此得到映射 $\hat G\to\mathrm{GL}_{N,\ZZ_\ell}$。于是引理说明它是闭浸入。现在我们有
\[
\begin{tikzcd}
\check G \arrow[rr] & & \mathrm{GL}_N \\
& \hat G \arrow[ul, dashed, "\exists"] \arrow[ur] &
\end{tikzcd}
\]
在 $\Ql$ 上，该映射延拓为 $\check G\to\hat G$；于是由平坦性，它也整地延拓。现在对 $\mathrm{SL}_2$ 运行与秩 $1$ 情形相同的论证：它在一般纤维上是同构，在 $\Fl$ 上是满射，因此在 $\ZZ_\ell$ 上是同构。\hfill$\square$
